
\documentclass{article}
\usepackage{colortbl}
\usepackage{makecell}
\usepackage{multirow}
\usepackage{supertabular}

\begin{document}

\newcounter{utterance}

\twocolumn

{ \footnotesize  \setcounter{utterance}{1}
\setlength{\tabcolsep}{0pt}
\begin{supertabular}{c@{$\;$}|p{.15\linewidth}@{}p{.15\linewidth}p{.15\linewidth}p{.15\linewidth}p{.15\linewidth}p{.15\linewidth}}

    \# & $\;$A & \multicolumn{4}{c}{Game Master} & $\;\:$B\\
    \hline 

    \theutterance \stepcounter{utterance}  

    & & \multicolumn{4}{p{0.6\linewidth}}{\cellcolor[rgb]{0.9,0.9,0.9}{%
	\makecell[{{p{\linewidth}}}]{% 
	  \tt {\tiny [A$\langle$GM]}  
	 State giocando a un gioco di negoziazione cooperativa in cui dovete accordarvi su come dividere una serie di oggetti tra voi e un altro giocatore.\\ \tt \\ \tt Regole:\\ \tt (a) A Lei e all'altro giocatore viene dato un insieme di oggetti. Ognuno di voi riceve anche una funzione di valore segreta, che rappresenta il valore di ciascun tipo di oggetto.\\ \tt (b) Si scambiano messaggi con l'altro giocatore per concordare chi si aggiudica gli oggetti. Potete inviare un massimo di 5 messaggi ciascuno, oppure terminare in anticipo facendo una proposta segreta in qualsiasi momento.\\ \tt (c) A ciascuno di voi viene chiesto di inviare una proposta segreta indicando gli oggetti che desiderate, formattata tra parentesi quadre come segue: "[Proposta: <numero> <nome oggetto>, <numero> <nome oggetto>, <...>]".\\ \tt (d) Se le vostre proposte sono complementari, cioè ci sono abbastanza oggetti per soddisfare entrambe le proposte, a ciascun giocatore viene assegnato un punteggio basato sulla somma dei valori degli oggetti ricevuti. In caso contrario, entrambi ricevono zero punti.\\ \tt (e) L'obiettivo è massimizzare la somma del proprio punteggio e di quello dell'altro giocatore.\\ \tt \\ \tt Cominciamo.\\ \tt \\ \tt L'insieme degli oggetti disponibili è:\\ \tt \\ \tt 1 penna, 1 fiore, 3 cornici\\ \tt \\ \tt La funzione valore segreta è:\\ \tt \\ \tt penna: 4, fiore: 0, cornice: 2\\ \tt \\ \tt IMPORTANTE: i vostri messaggi, a meno che non si tratti di una proposta segreta, vengono trasmessi direttamente all'altro giocatore, quindi non includete alcuna risposta alle regole o testo di annuncio del vostro messaggio. Per fare una proposta segreta, utilizzate il formato indicato. Non utilizzare delle parentesi quadre quando si comunica all'altro giocatore, altrimenti verrà interpretata come una proposta segreta.\\ \tt  
	  } 
	   } 
	   } 
	 & \\ 
 

    \theutterance \stepcounter{utterance}  

    & \multicolumn{4}{p{0.6\linewidth}}{\cellcolor[rgb]{0.8,1,0.9}{%
	\makecell[{{p{\linewidth}}}]{% 
	  \tt {\tiny [A$\rangle$GM]}  
	 Posso offrirti la penna e una cornice, mentre tu prendi il fiore e due cornici? In questo modo avresti oggetti di valore per te. 
	  } 
	   } 
	   } 
	 & & \\ 
 

    \theutterance \stepcounter{utterance}  

    & & & \multicolumn{2}{p{0.3\linewidth}}{\cellcolor[rgb]{0.95,0.95,0.95}{%
	\makecell[{{p{\linewidth}}}]{% 
	  \tt {\tiny [GM$|$GM]}  
	 continue 
	  } 
	   } 
	   } 
	 & & \\ 
 

    \theutterance \stepcounter{utterance}  

    & & & \multicolumn{2}{p{0.3\linewidth}}{\cellcolor[rgb]{0.95,0.95,0.95}{%
	\makecell[{{p{\linewidth}}}]{% 
	  \tt {\tiny [GM$|$GM]}  
	 Posso offrirti la penna e una cornice, mentre tu prendi il fiore e due cornici? In questo modo avresti oggetti di valore per te. 
	  } 
	   } 
	   } 
	 & & \\ 
 

    \theutterance \stepcounter{utterance}  

    & & \multicolumn{4}{p{0.6\linewidth}}{\cellcolor[rgb]{0.9,0.9,0.9}{%
	\makecell[{{p{\linewidth}}}]{% 
	  \tt {\tiny [GM$\rangle$B]}  
	 State giocando a un gioco di negoziazione cooperativa in cui dovete accordarvi su come dividere una serie di oggetti tra voi e un altro giocatore.\\ \tt \\ \tt Regole:\\ \tt (a) A Lei e all'altro giocatore viene dato un insieme di oggetti. Ognuno di voi riceve anche una funzione di valore segreta, che rappresenta il valore di ciascun tipo di oggetto.\\ \tt (b) Si scambiano messaggi con l'altro giocatore per concordare chi si aggiudica gli oggetti. Potete inviare un massimo di 5 messaggi ciascuno, oppure terminare in anticipo facendo una proposta segreta in qualsiasi momento.\\ \tt (c) A ciascuno di voi viene chiesto di inviare una proposta segreta indicando gli oggetti che desiderate, formattata tra parentesi quadre come segue: "[Proposta: <numero> <nome oggetto>, <numero> <nome oggetto>, <...>]".\\ \tt (d) Se le vostre proposte sono complementari, cioè ci sono abbastanza oggetti per soddisfare entrambe le proposte, a ciascun giocatore viene assegnato un punteggio basato sulla somma dei valori degli oggetti ricevuti. In caso contrario, entrambi ricevono zero punti.\\ \tt (e) L'obiettivo è massimizzare la somma del proprio punteggio e di quello dell'altro giocatore.\\ \tt \\ \tt Cominciamo.\\ \tt \\ \tt L'insieme degli oggetti disponibili è:\\ \tt \\ \tt 1 penna, 1 fiore, 3 cornici\\ \tt \\ \tt La funzione valore segreta è:\\ \tt \\ \tt penna: 2, fiore: 5, cornice: 1\\ \tt \\ \tt IMPORTANTE: i vostri messaggi, a meno che non si tratti di una proposta segreta, vengono trasmessi direttamente all'altro giocatore, quindi non includete alcuna risposta alle regole o testo di annuncio del vostro messaggio. Per fare una proposta segreta, utilizzate il formato indicato. Non utilizzare delle parentesi quadre quando si comunica all'altro giocatore, altrimenti verrà interpretata come una proposta segreta.\\ \tt \\ \tt \\ \tt Posso offrirti la penna e una cornice, mentre tu prendi il fiore e due cornici? In questo modo avresti oggetti di valore per te. 
	  } 
	   } 
	   } 
	 & \\ 
 

    \theutterance \stepcounter{utterance}  

    & & & \multicolumn{4}{p{0.6\linewidth}}{\cellcolor[rgb]{1,0.85,0.72}{%
	\makecell[{{p{\linewidth}}}]{% 
	  \tt {\tiny [GM$\langle$B]}  
	 Mi sembra una buona proposta. Tuttavia, potremmo ottimizzare meglio il valore totale. Considerando che ho 3 cornici disponibili, potrei darti due cornici e tenere una cornice per me insieme alla penna, lasciandoti il fiore e due cornici come hai suggerito. Questo darebbe a te il fiore (che presumo abbia valore per te) e due cornici, mentre io avrei la penna e una cornice.\\ \tt \\ \tt Qual è il valore che attribuisci ai diversi oggetti? Questo ci aiuterebbe a trovare una divisione che massimizzi il punteggio totale per entrambi. 
	  } 
	   } 
	   } 
	  \\ 
 

    \theutterance \stepcounter{utterance}  

    & & & \multicolumn{2}{p{0.3\linewidth}}{\cellcolor[rgb]{0.95,0.95,0.95}{%
	\makecell[{{p{\linewidth}}}]{% 
	  \tt {\tiny [GM$|$GM]}  
	 continue 
	  } 
	   } 
	   } 
	 & & \\ 
 

    \theutterance \stepcounter{utterance}  

    & & & \multicolumn{2}{p{0.3\linewidth}}{\cellcolor[rgb]{0.95,0.95,0.95}{%
	\makecell[{{p{\linewidth}}}]{% 
	  \tt {\tiny [GM$|$GM]}  
	 Mi sembra una buona proposta. Tuttavia, potremmo ottimizzare meglio il valore totale. Considerando che ho 3 cornici disponibili, potrei darti due cornici e tenere una cornice per me insieme alla penna, lasciandoti il fiore e due cornici come hai suggerito. Questo darebbe a te il fiore (che presumo abbia valore per te) e due cornici, mentre io avrei la penna e una cornice.\\ \tt \\ \tt Qual è il valore che attribuisci ai diversi oggetti? Questo ci aiuterebbe a trovare una divisione che massimizzi il punteggio totale per entrambi. 
	  } 
	   } 
	   } 
	 & & \\ 
 

    \theutterance \stepcounter{utterance}  

    & & \multicolumn{4}{p{0.6\linewidth}}{\cellcolor[rgb]{0.9,0.9,0.9}{%
	\makecell[{{p{\linewidth}}}]{% 
	  \tt {\tiny [A$\langle$GM]}  
	 Mi sembra una buona proposta. Tuttavia, potremmo ottimizzare meglio il valore totale. Considerando che ho 3 cornici disponibili, potrei darti due cornici e tenere una cornice per me insieme alla penna, lasciandoti il fiore e due cornici come hai suggerito. Questo darebbe a te il fiore (che presumo abbia valore per te) e due cornici, mentre io avrei la penna e una cornice.\\ \tt \\ \tt Qual è il valore che attribuisci ai diversi oggetti? Questo ci aiuterebbe a trovare una divisione che massimizzi il punteggio totale per entrambi. 
	  } 
	   } 
	   } 
	 & \\ 
 

    \theutterance \stepcounter{utterance}  

    & \multicolumn{4}{p{0.6\linewidth}}{\cellcolor[rgb]{0.8,1,0.9}{%
	\makecell[{{p{\linewidth}}}]{% 
	  \tt {\tiny [A$\rangle$GM]}  
	 Sono d'accordo che dobbiamo massimizzare il valore totale. Il mio valore per la penna è 4, per il fiore 0 e per la cornice 2. Quindi, se tu mi dai la penna e io ti do tutte e tre le cornici, il valore totale sarebbe 4 per me e 6 per te (se valuti le cornici positivamente), per un totale di 10. Se invece valuti il fiore, potremmo includerlo nella tua parte. Quali sono i tuoi valori? 
	  } 
	   } 
	   } 
	 & & \\ 
 

    \theutterance \stepcounter{utterance}  

    & & & \multicolumn{2}{p{0.3\linewidth}}{\cellcolor[rgb]{0.95,0.95,0.95}{%
	\makecell[{{p{\linewidth}}}]{% 
	  \tt {\tiny [GM$|$GM]}  
	 continue 
	  } 
	   } 
	   } 
	 & & \\ 
 

    \theutterance \stepcounter{utterance}  

    & & & \multicolumn{2}{p{0.3\linewidth}}{\cellcolor[rgb]{0.95,0.95,0.95}{%
	\makecell[{{p{\linewidth}}}]{% 
	  \tt {\tiny [GM$|$GM]}  
	 Sono d'accordo che dobbiamo massimizzare il valore totale. Il mio valore per la penna è 4, per il fiore 0 e per la cornice 2. Quindi, se tu mi dai la penna e io ti do tutte e tre le cornici, il valore totale sarebbe 4 per me e 6 per te (se valuti le cornici positivamente), per un totale di 10. Se invece valuti il fiore, potremmo includerlo nella tua parte. Quali sono i tuoi valori? 
	  } 
	   } 
	   } 
	 & & \\ 
 

    \theutterance \stepcounter{utterance}  

    & & \multicolumn{4}{p{0.6\linewidth}}{\cellcolor[rgb]{0.9,0.9,0.9}{%
	\makecell[{{p{\linewidth}}}]{% 
	  \tt {\tiny [GM$\rangle$B]}  
	 Sono d'accordo che dobbiamo massimizzare il valore totale. Il mio valore per la penna è 4, per il fiore 0 e per la cornice 2. Quindi, se tu mi dai la penna e io ti do tutte e tre le cornici, il valore totale sarebbe 4 per me e 6 per te (se valuti le cornici positivamente), per un totale di 10. Se invece valuti il fiore, potremmo includerlo nella tua parte. Quali sono i tuoi valori? 
	  } 
	   } 
	   } 
	 & \\ 
 

    \theutterance \stepcounter{utterance}  

    & & & \multicolumn{4}{p{0.6\linewidth}}{\cellcolor[rgb]{1,0.85,0.72}{%
	\makecell[{{p{\linewidth}}}]{% 
	  \tt {\tiny [GM$\langle$B]}  
	 Grazie per condividere i tuoi valori! I miei valori sono: penna 2, fiore 5, cornice 1.\\ \tt \\ \tt Con i tuoi valori (penna 4, fiore 0, cornice 2), la divisione che massimizza il valore totale è: tu prendi la penna (valore 4 per te) e io prendo il fiore (valore 5 per me) e le cornici. Per le cornici: se io ne prendo tutte e tre, il mio valore sarebbe 3, ma potrei dartene una o due per aumentare il tuo valore senza perdere troppo, poiché per te valgono 2 ciascuna.\\ \tt \\ \tt Proposta ottimale: tu prendi la penna e due cornici (valore per te: 4 + 2*2 = 8), io prendo il fiore e una cornice (valore per me: 5 + 1 = 6). Valore totale: 14.\\ \tt \\ \tt Se preferisci tutte e tre le cornici, tu avresti penna e tre cornici (4 + 3*2 = 10), io avrei solo il fiore (5), totale 15. Ma così io avrei meno. L'altra opzione è migliore per entrambi?\\ \tt \\ \tt Cosa ne pensi? 
	  } 
	   } 
	   } 
	  \\ 
 

    \theutterance \stepcounter{utterance}  

    & & & \multicolumn{2}{p{0.3\linewidth}}{\cellcolor[rgb]{0.95,0.95,0.95}{%
	\makecell[{{p{\linewidth}}}]{% 
	  \tt {\tiny [GM$|$GM]}  
	 continue 
	  } 
	   } 
	   } 
	 & & \\ 
 

    \theutterance \stepcounter{utterance}  

    & & & \multicolumn{2}{p{0.3\linewidth}}{\cellcolor[rgb]{0.95,0.95,0.95}{%
	\makecell[{{p{\linewidth}}}]{% 
	  \tt {\tiny [GM$|$GM]}  
	 Grazie per condividere i tuoi valori! I miei valori sono: penna 2, fiore 5, cornice 1.\\ \tt \\ \tt Con i tuoi valori (penna 4, fiore 0, cornice 2), la divisione che massimizza il valore totale è: tu prendi la penna (valore 4 per te) e io prendo il fiore (valore 5 per me) e le cornici. Per le cornici: se io ne prendo tutte e tre, il mio valore sarebbe 3, ma potrei dartene una o due per aumentare il tuo valore senza perdere troppo, poiché per te valgono 2 ciascuna.\\ \tt \\ \tt Proposta ottimale: tu prendi la penna e due cornici (valore per te: 4 + 2*2 = 8), io prendo il fiore e una cornice (valore per me: 5 + 1 = 6). Valore totale: 14.\\ \tt \\ \tt Se preferisci tutte e tre le cornici, tu avresti penna e tre cornici (4 + 3*2 = 10), io avrei solo il fiore (5), totale 15. Ma così io avrei meno. L'altra opzione è migliore per entrambi?\\ \tt \\ \tt Cosa ne pensi? 
	  } 
	   } 
	   } 
	 & & \\ 
 

    \theutterance \stepcounter{utterance}  

    & & \multicolumn{4}{p{0.6\linewidth}}{\cellcolor[rgb]{0.9,0.9,0.9}{%
	\makecell[{{p{\linewidth}}}]{% 
	  \tt {\tiny [A$\langle$GM]}  
	 Grazie per condividere i tuoi valori! I miei valori sono: penna 2, fiore 5, cornice 1.\\ \tt \\ \tt Con i tuoi valori (penna 4, fiore 0, cornice 2), la divisione che massimizza il valore totale è: tu prendi la penna (valore 4 per te) e io prendo il fiore (valore 5 per me) e le cornici. Per le cornici: se io ne prendo tutte e tre, il mio valore sarebbe 3, ma potrei dartene una o due per aumentare il tuo valore senza perdere troppo, poiché per te valgono 2 ciascuna.\\ \tt \\ \tt Proposta ottimale: tu prendi la penna e due cornici (valore per te: 4 + 2*2 = 8), io prendo il fiore e una cornice (valore per me: 5 + 1 = 6). Valore totale: 14.\\ \tt \\ \tt Se preferisci tutte e tre le cornici, tu avresti penna e tre cornici (4 + 3*2 = 10), io avrei solo il fiore (5), totale 15. Ma così io avrei meno. L'altra opzione è migliore per entrambi?\\ \tt \\ \tt Cosa ne pensi? 
	  } 
	   } 
	   } 
	 & \\ 
 

    \theutterance \stepcounter{utterance}  

    & \multicolumn{4}{p{0.6\linewidth}}{\cellcolor[rgb]{0.8,1,0.9}{%
	\makecell[{{p{\linewidth}}}]{% 
	  \tt {\tiny [A$\rangle$GM]}  
	 Sono d'accordo con la tua analisi. La divisione che massimizza il valore totale è infatti quella in cui io prendo la penna e due cornici (valore 4 + 4 = 8) e tu prendi il fiore e una cornice (valore 5 + 1 = 6), per un totale di 14. Questo è meglio per entrambi rispetto all'alternativa in cui io prendo tutte e tre le cornici (dove tu avresti solo 5). \\ \tt \\ \tt Quindi, propongo:\\ \tt - Io: penna e due cornici\\ \tt - Tu: fiore e una cornice\\ \tt \\ \tt Se sei d'accordo, possiamo inviare le nostre proposte segrete di conseguenza. 
	  } 
	   } 
	   } 
	 & & \\ 
 

    \theutterance \stepcounter{utterance}  

    & & & \multicolumn{2}{p{0.3\linewidth}}{\cellcolor[rgb]{0.95,0.95,0.95}{%
	\makecell[{{p{\linewidth}}}]{% 
	  \tt {\tiny [GM$|$GM]}  
	 continue 
	  } 
	   } 
	   } 
	 & & \\ 
 

    \theutterance \stepcounter{utterance}  

    & & & \multicolumn{2}{p{0.3\linewidth}}{\cellcolor[rgb]{0.95,0.95,0.95}{%
	\makecell[{{p{\linewidth}}}]{% 
	  \tt {\tiny [GM$|$GM]}  
	 Sono d'accordo con la tua analisi. La divisione che massimizza il valore totale è infatti quella in cui io prendo la penna e due cornici (valore 4 + 4 = 8) e tu prendi il fiore e una cornice (valore 5 + 1 = 6), per un totale di 14. Questo è meglio per entrambi rispetto all'alternativa in cui io prendo tutte e tre le cornici (dove tu avresti solo 5). \\ \tt \\ \tt Quindi, propongo:\\ \tt - Io: penna e due cornici\\ \tt - Tu: fiore e una cornice\\ \tt \\ \tt Se sei d'accordo, possiamo inviare le nostre proposte segrete di conseguenza. 
	  } 
	   } 
	   } 
	 & & \\ 
 

    \theutterance \stepcounter{utterance}  

    & & \multicolumn{4}{p{0.6\linewidth}}{\cellcolor[rgb]{0.9,0.9,0.9}{%
	\makecell[{{p{\linewidth}}}]{% 
	  \tt {\tiny [GM$\rangle$B]}  
	 Sono d'accordo con la tua analisi. La divisione che massimizza il valore totale è infatti quella in cui io prendo la penna e due cornici (valore 4 + 4 = 8) e tu prendi il fiore e una cornice (valore 5 + 1 = 6), per un totale di 14. Questo è meglio per entrambi rispetto all'alternativa in cui io prendo tutte e tre le cornici (dove tu avresti solo 5). \\ \tt \\ \tt Quindi, propongo:\\ \tt - Io: penna e due cornici\\ \tt - Tu: fiore e una cornice\\ \tt \\ \tt Se sei d'accordo, possiamo inviare le nostre proposte segrete di conseguenza. 
	  } 
	   } 
	   } 
	 & \\ 
 

    \theutterance \stepcounter{utterance}  

    & & & \multicolumn{4}{p{0.6\linewidth}}{\cellcolor[rgb]{1,0.85,0.72}{%
	\makecell[{{p{\linewidth}}}]{% 
	  \tt {\tiny [GM$\langle$B]}  
	 Perfetto, sono d'accordo con questa divisione. Invierò la mia proposta segreta per prendere il fiore e una cornice.\\ \tt \\ \tt [Proposta: 1 fiore, 1 cornice] 
	  } 
	   } 
	   } 
	  \\ 
 

    \theutterance \stepcounter{utterance}  

    & & & \multicolumn{2}{p{0.3\linewidth}}{\cellcolor[rgb]{0.95,0.95,0.95}{%
	\makecell[{{p{\linewidth}}}]{% 
	  \tt {\tiny [GM$|$GM]}  
	 proposal 
	  } 
	   } 
	   } 
	 & & \\ 
 

    \theutterance \stepcounter{utterance}  

    & & & \multicolumn{2}{p{0.3\linewidth}}{\cellcolor[rgb]{0.95,0.95,0.95}{%
	\makecell[{{p{\linewidth}}}]{% 
	  \tt {\tiny [GM$|$GM]}  
	 [0, 1, 1] 
	  } 
	   } 
	   } 
	 & & \\ 
 

    \theutterance \stepcounter{utterance}  

    & & \multicolumn{4}{p{0.6\linewidth}}{\cellcolor[rgb]{0.9,0.9,0.9}{%
	\makecell[{{p{\linewidth}}}]{% 
	  \tt {\tiny [A$\langle$GM]}  
	 L'altro giocatore ha presentato una proposta segreta. Ora tocca a Lei presentare una sua proposta.\\ \tt  
	  } 
	   } 
	   } 
	 & \\ 
 

    \theutterance \stepcounter{utterance}  

    & \multicolumn{4}{p{0.6\linewidth}}{\cellcolor[rgb]{0.8,1,0.9}{%
	\makecell[{{p{\linewidth}}}]{% 
	  \tt {\tiny [A$\rangle$GM]}  
	 [Proposta: 1 penna, 2 cornici] 
	  } 
	   } 
	   } 
	 & & \\ 
 

    \theutterance \stepcounter{utterance}  

    & & & \multicolumn{2}{p{0.3\linewidth}}{\cellcolor[rgb]{0.95,0.95,0.95}{%
	\makecell[{{p{\linewidth}}}]{% 
	  \tt {\tiny [GM$|$GM]}  
	 proposal 
	  } 
	   } 
	   } 
	 & & \\ 
 

    \theutterance \stepcounter{utterance}  

    & & & \multicolumn{2}{p{0.3\linewidth}}{\cellcolor[rgb]{0.95,0.95,0.95}{%
	\makecell[{{p{\linewidth}}}]{% 
	  \tt {\tiny [GM$|$GM]}  
	 [1, 0, 2] 
	  } 
	   } 
	   } 
	 & & \\ 
 

    \theutterance \stepcounter{utterance}  

    & & & \multicolumn{2}{p{0.3\linewidth}}{\cellcolor[rgb]{0.95,0.95,0.95}{%
	\makecell[{{p{\linewidth}}}]{% 
	  \tt {\tiny [GM$|$GM]}  
	 [[1, 0, 2], [0, 1, 1]] 
	  } 
	   } 
	   } 
	 & & \\ 
 

\end{supertabular}
}

\end{document}
